% Source: sources/returns/S6_04_MAPS_MANIFOLDS/sections/manifolds.tex:6823–7060
\paragraph{Theorem D.2.33 (Proved: all original multiplicative choices on the natural catalogue).}
At every zero-free multiplicative site in the $172$ natural
initial rows, either the original cap or the marked primitive
Mumford cap may be used. Together with all the original and
finite logarithmic additive choices of D.2.32, retaining
original caps at $I_n^*$, this gives $495$ distinct integral homology-presentation
patterns, after identifying the boundary-kernel agreements
for zero-degree original/Mumford choices proved below.
Every original/Mumford choice represented by such an
agreement is also listed explicitly in
\path{records/manifolds_original_multiplicative.json}.
Every finite and regular logarithmic order and every free
affine twist remains allowed.

The new original caps occur at $120$ marked zero-free
$I_n$ sites with nonzero component Chern vector. Each adds
\[
 h_i s=0,\qquad h_i=b/n,                                   \tag{D.OM1}
\]
to the exact global matrix, has free local degree-two
cokernel of rank $n-2$, and adds one to the free
degree-one Mayer--Vietoris bridge. Consequently D.2.32
continues to hold verbatim if $L$ counts \emph{all}
nonzero-Chern original caps, additive or multiplicative,
and $\beta$ includes $n-2$ for each new multiplicative
original. In particular, when $L>0$,
\[
\begin{array}{c|c}
j&H_j(X;\mathbb Z)\\\hline
0,6&\mathbb Z\\
1&\mathbb Z^r\oplus A_1\\
2&\mathbb Z^{\beta+L+r}\oplus T\\
3&\mathbb Z^{2\beta+2L}\oplus T\\
4&\mathbb Z^{\beta+L+r}\oplus A_1\\
5&\mathbb Z^r.
\end{array}                                                \tag{D.OM2}
\]
Here $\mathbb Z^r\oplus A_1$ is the exact Seifert
abelianization of D.2.23, and $T$ is the torsion of the
full finite/regular normal matrix with all the additive
relations (D.OC1) and multiplicative relations (D.OM1).
The full fundamental group is still the Seifert group,
including when its rational Euler number is zero.

\begin{proof}
\emph{The actual bundle over the nodal cycle.}
At a zero-free $I_n$ site in these rows, either $P$ is
narrow or $Q$ is narrow. If $P$ is narrow, $Q$ acts
transitively on the components. Its original cap has
one component orbit and the same full integral normal
kernel and primitive fibre image as the marked Mumford
cap, by the multiplicative calculation in D.2.31.
This proves equality of the recorded integral homology
presentations for these zero-degree choices. It does not
assert a relative diffeomorphism between their caps.

Otherwise $n>1$, $Q$ is narrow and the component of $P$
is a unit modulo $n$. Write $K=I_n$ and
\[
 c=e_{\overline P}-e_0\in H^2(K;\mathbb Z)=\mathbb Z^n .
\]
This vector is primitive and its sum is zero. The local
Neron logarithm of $Q$ contracts its translation to the
identity. The lift preserves the line bundle with Chern
vector $c$; this bundle is not replaced by a trivial one.
There is one additional point compared with the tree
argument of D.2.32: the dual graph of $I_n$ has a circle.
On a cycle, translation in the identity component acts
on the generalized Jacobian of a multidegree-$c$ line
by the power given by $\sum_i c_i$. That sum is zero,
so the translation isotopy admits line-bundle lifts.
One can see this directly on the normalization: multiply
the coordinate on every rational component by the same
unit and use its integral powers in the line transition
functions; the final clutch factor is the unit to
$\sum_i c_i=0$.

The ratio between the resulting lift and the prescribed
holomorphic lift on $K$ is a global holomorphic unit on
the compact connected nodal cycle, and hence is constant.
It has no winding on the dual circle. Rigidification
removes this constant phase, and the nonvanishing
linearization on the base disc has a logarithm, fixing
the normal-meridian marking as well. Thus no additional
mixed Chern class is introduced. The actual cap retracts,
with its transported boundary marking, to
\[
 U\simeq Y_c\times S^1_\gamma,\qquad
 S^1_\delta\longrightarrow Y_c\longrightarrow K,           \tag{D.OM3}
\]
where the circle bundle has Chern class $c$.

The integral Gysin sequence, or the bundle spectral
sequence over $K\simeq S^1\vee\bigvee_{i=1}^nS^2$, gives
\[
\begin{aligned}
 H_1(Y_c)&=\mathbb Z,&
 H_2(Y_c)&=\mathbb Z^n,&
 H_3(Y_c)&=\mathbb Z^n,\\
 H_1(U)&=\mathbb Z^2,&
 H_2(U)&=\mathbb Z^{n+1},&
 H_3(U)&=\mathbb Z^{2n},&
 H_4(U)&=\mathbb Z^n .
\end{aligned}                                             \tag{D.OM4}
\]
All groups are free: the only transgression is the
primitive vector $c$. More precisely,
\[
 0\longrightarrow\mathbb Z[\delta\times\text{dual circle}]
 \longrightarrow H_2(Y_c)
 \longrightarrow\ker(c:\mathbb Z^n\to\mathbb Z)
 \longrightarrow0                                        \tag{D.OM5}
\]
is an exact sequence of free groups. This also identifies
its first generator primitively.

\emph{The full marked kernel, including its normal correction.}
Let $u_1$ be the primitive elliptic vanishing circle and
choose $u_2$ with oriented determinant one. As in D.2.32,
narrowness gives an integral vector $z$ satisfying
$(B-I)z=Q_{\rm col}$, and put $\ell=(-z,0,1)$.
In the integral basis $(u_1,u_2,\delta,\ell)$ the actual
monodromy has the form
\[
 A u_2=u_2+\nu u_1+p\delta,\qquad
 A u_1=u_1,\quad A\delta=\delta,\quad A\ell=\ell,
 \qquad |\nu|=n,\quad\gcd(\nu,p)=1.                        \tag{D.OM6}
\]
These are equalities in the marked lattice, not only
a rational normal form. The certificate records the
unimodular basis for each of the $120$ occurrences.
In particular the full fixed lattice is
$\mathbb Zu_1\oplus\mathbb Z\delta\oplus\mathbb Z\ell$,
and both $A-I$ and $\bigwedge^2 A-I$ have only unit
nonzero Smith factors.

The regular fibre has primitive rank-three image in
$H_2(U)$. Its generators can be taken as
\[
 u_1\wedge u_2,\qquad u_2\wedge\delta,\qquad
 u_2\wedge\ell.                                           \tag{D.OM7}
\]
The first projects to the all-one component vector in
$\ker c$, which is primitive; a possible addition of
the first term of (D.OM5) does not affect this assertion.
The second generates that first term primitively, and
the third is the dual-circle class times $\gamma$.
Hence their span is primitive. The full regular-fibre
kernel is exactly
\[
 K_{\rm fib}=
 \mathbb Z(u_1\wedge\delta)\oplus
 \mathbb Z(u_1\wedge\ell)\oplus
 \mathbb Z(\delta\wedge\ell).                              \tag{D.OM8}
\]
For the first two generators, the vanishing disc and
its lifted products provide the required boundaries.
For the third, primitivity of $c$ kills the circle
$\delta$ integrally in $Y_c$. The image of
$\bigwedge^2 A-I$ is a primitive rank-two sublattice
of (D.OM8).

The normal tori in the $\delta$ and $\ell$ directions
bound the sweeps of the origin disc under the two
principal circle actions. Their marked fibre corrections
are zero. A normal torus in the $u_1$ direction requires
slightly more care. Its projection to the original
elliptic neighbourhood is the usual vanishing-loop
normal torus, which bounds by the local vanishing-disc
construction. Its image in $H_2(U)$ therefore lies in
the first summand of (D.OM5); its $\gamma$ coordinate
is zero. It is an integral multiple of the regular
torus $u_2\wedge\delta$. Subtracting that multiple
produces an actual boundary-kernel cycle with normal
projection $u_1$.

This correction cannot be discarded at the local level.
However the global primitive alternating functional
satisfies exactly
\[
 \omega(u_2,\delta)=0,\qquad
 \omega=a e_1^*e_2^*-b\delta^*\gamma^*.                    \tag{D.OM9}
\]
Both $u_2$ and $\delta$ have zero $\gamma$ coordinate,
and $\delta$ has zero elliptic projection. Thus the
correction has zero image in the exterior fibre group.
The actual global image of the three normal kernels is
therefore the full marked fixed-sweep lattice with
zero fibre coefficient. This proves the needed marked
statement without guessing the local correction integer.

The boundary image is exactly the primitive rank-three
regular-fibre image (D.OM7), and consequently
\[
 H_2(U)/\operatorname{im}H_2(\partial U)
                          =\mathbb Z^{n-2}.               \tag{D.OM10}
\]
Evaluate (D.OM8) by $\omega$. Monodromy invariance gives
$\nu\,\omega(u_1,\ell)+p\,\omega(\delta,\ell)=0$,
where $\omega(\delta,\ell)=-b$ and
$\omega(u_1,\delta)=0$. Since $\gcd(\nu,p)=1$,
$n\mid b$ and the exact common divisor of these values
is $b/n$. This proves (D.OM1), including its integral
coefficient.

\emph{Global gluing and exhaustion of the stated choices.}
The cap has first Betti number two, whereas the fixed
lattice in (D.OM6) has rank three. Thus it adds precisely
one free bridge, just as a nonzero-Chern additive original
does. The local degree-two cokernel (D.OM10) is free,
so it splits from the remaining global presentation.
The preceding kernel calculation shows that the latter
is exactly the old fixed-sweep matrix with the new
relation (D.OM1). All finite torsion caps and their
nontrivial normal Chern classes are still handled by
the full three-column Thom matrices of D.2.30.

The proof of (D.OC3) now applies with the enlarged $L$
and $\beta$. In particular, a positive $L$ kills a
nonzero multiple of $s$, so the reduced degree-two
matrix has rational cokernel zero, even when $r=1$.
The degree-one bridge has rank $L+r$. Every original
cap in (D.OM3) has Euler number zero, and the unique
index-one toric zero still contributes two. Integral
duality therefore gives (D.OM2).
At the group level the new original cap only adds
vertical relations, including the vanishing direction
and $\delta$, and adds no relation on $\gamma$.
The complete vertical lattice was already killed in
D.2.23, so its Seifert presentation is unchanged.

Finally, choose original or marked Mumford independently
at each of the nonzero-Chern multiplicative sites in
the $282$ additive patterns. This yields $495$ integral homology-presentation patterns
on all $172$ initial rows, with all the zero-degree
original/Mumford alternatives separately recorded. It
exhausts exactly the choices stated in the theorem.
\end{proof}
